\chapter{A interface}

\eyebrow{a interface}

\vspace{4pt}
O Apex segue a anatomia do software de controle de missão. A janela é organizada em
regiões fixas, e a operação consiste em \textbf{selecionar um objeto} na árvore de
navegação e \textbf{escolher uma vista} dele.

\section{Anatomia da janela}

\begin{center}
\begin{tabularx}{\linewidth}{@{}l X@{}}
\toprule
\textbf{Região} & \textbf{O que ela contém} \\
\midrule
Barra de título (topo) & Marca do produto e uma leitura ao vivo: carcaças, imagens, avaliadas. \\
Árvore de navegação (esquerda) & Objetos raiz (Monitor ao vivo, Importar imagens, Painel do projeto) e a hierarquia de dados \textbf{Lotes\,$\rightarrow$\,Carcaças}. \\
Barra de navegação (acima do centro) & O caminho (trilha) do objeto selecionado, e um \textbf{seletor de vistas} à direita. \\
Área principal (centro) & A vista ativa do objeto selecionado. \\
Inspetor (direita) & Propriedades somente leitura do objeto selecionado. Recolhível. \\
Barra de status (inferior) & Estado do motor, dispositivo e contagens do conjunto de dados. \\
\bottomrule
\end{tabularx}
\end{center}

\section{Objetos e vistas}

Selecionar um objeto na árvore altera tanto a trilha quanto o conjunto de vistas
oferecidas na barra de navegação:

\begin{center}
\begin{tabularx}{\linewidth}{@{}l X@{}}
\toprule
\textbf{Objeto} & \textbf{Vistas (na barra de navegação)} \\
\midrule
Um \textbf{lote}      & Visão geral, Análise, Avaliação \\
Uma \textbf{carcaça}  & Visão geral, Imagens, Análise \\
Monitor ao vivo       & Monitor ao vivo \\
Importar imagens      & Importação \\
Painel do projeto     & Painel \\
\bottomrule
\end{tabularx}
\end{center}

\begin{note}[title={O Inspetor acompanha o objeto selecionado}]
Qualquer que seja o objeto selecionado, um lote ou uma carcaça, o Inspetor à
direita mostra suas propriedades: para uma carcaça, sua identificação, referência
física e análise de modelo; para um lote, suas contagens e estatísticas de gordura.
Ele é somente leitura; a edição ocorre nas vistas.
\end{note}

A barra de status na parte inferior mostra o estado corrente do instrumento: indica
se o \tele{engine} de inferência está \pillok{ready} e em qual \tele{device} (por
exemplo, \tele{mps}), além dos totais correntes do conjunto de dados.
